% ---------------------------------------------------------------------
\section{证明策略}\label{sec:method-strategy}

因子化陈述（\cref{def:intro-factorization}）在微扰论中应逐阶地对红外正规化的部分子态成立。对固定阶的一张图（或具有公共圈动量走向的一族图），它意味着被积函数 $\Gamma$
可以写成若干减除项之和加一个余项，余项的积分在 $\lambda\to0$ 时被幂次压低、且对 $\varepsilon$ 一致（误差类型 E-N，\cref{def:intro-error-types}），而减除项由 Ward 恒等式化为硬部分与算符矩阵元的卷积。
于是证明分成两件事：\emph{幂次计数与减除}（余项小）与\emph{规范结构}（减除项可分解）。传统的全阶论证（Libby--Sterman\cite{LibbySterman1978a}、Collins--Soper--Sterman\cite{CollinsSoperSterman1989}、
Collins\cite{Collins2011}）可归纳为六步（\cref{fig:method-flowchart}）：
\begin{enumerate}[label=(\arabic*),leftmargin=2.2em,itemsep=0pt,topsep=2pt]
\item \emph{区域}：由 Landau 方程与幂次计数（\cref{prop:ir-power-counting}）找出领头区域（\cref{thm:ir-leading-regions}），本章把它们表述为逐线的模式指派（\cref{sec:method-regions}）；
\item \emph{近似}：用幂等投影作共线近似与软近似；近似子不引入新分母，被近似的传播子保留原有的 $i0$——它决定此后 Wilson 线的方向；
\item \emph{Grammer--Yennie 分解}：连接共线（或软）子图与硬部分的胶子分为 K、G 两部分，K 部分适用 Ward 恒等式，G 部分被压低（\cref{sec:method-gy}）；
\item \emph{Ward 恒等式}：K 胶子从硬部分退耦，改为附着在 Wilson 线上（\cref{sec:method-ward}）；次序是“先近似、后用 Ward 恒等式、再得 Wilson 线”，
  在近似之前对原图使用 Ward 恒等式会产生不受控制的项\cite{CollinsSoperSterman1989,Labastida:1984gy}；
\item \emph{Glauber 区、幺正性与减除}：把围道移出 Glauber 区或证明其贡献相消（\cref{sec:method-glauber}），末态相互作用由幺正性相消（\cref{sec:method-unitarity}），
  区域重叠的重复计数由 Collins 的逐级减除消除（\cref{sec:method-subtraction}）；
\item \emph{因子化形式}：对图与区域求和，得到 $\sigma=\sum_{ij}H_{ij}\otimes f_{i/A}\otimes f_{j/B}\,[\otimes S]+\order{\lambda^p}$（\cref{sec:method-template}）。
\end{enumerate}
第 (5) 步中的次序是固定的：软因子的幺正性（它对共线单举截面给出软因子为 $1$）只能在处理完 Glauber 区之后使用；同样，DGLAP 方程来自算符重整化（\cref{thm:pdf-dglap}），
因子化标度无关性是其推论（\cref{prop:method-scale-scheme}）。

每一步都在两个层次上陈述：文献中的全阶论证层次——幂次计数给出的是\emph{标度陈述}（E-S），结论为 \tagA；以及精确的条件层次——把“余项小”写成积分范数下的界（E-N），
把不能证明的几何—解析输入写成编号假设，再给出完整的演绎（\tagR，$R_{\mathrm{cond}}$，见附录~N 与附录~O）。两个层次之间的差距恰好就是\cref{asm:appN-envelope,asm:appN-separation,asm:appN-collins}的内容；
章末的\cref{tab:method-correspondence}列出二者的对应。

